\documentclass[11pt]{article}
\usepackage[margin=2.3cm]{geometry}
\usepackage{amssymb}
\usepackage[dvipsnames,table]{xcolor}
\usepackage{enumitem}
\usepackage{titlesec}
\usepackage{parskip}
\usepackage{fancyhdr}
\usepackage{hyperref}
\hypersetup{colorlinks=true, linkcolor=RoyalBlue, urlcolor=RoyalBlue}
\definecolor{deep}{HTML}{1F4E79}
\definecolor{soft}{HTML}{EAF1FA}
\definecolor{warn}{HTML}{FBE9E7}
\definecolor{tipc}{HTML}{E8F5E9}
\titleformat{\section}{\Large\bfseries\color{deep}}{\thesection.}{0.5em}{}
\titleformat{\subsection}{\large\bfseries\color{deep}}{\thesubsection}{0.5em}{}
\titleformat{\subsubsection}{\normalsize\bfseries}{}{0em}{}
\setlist{leftmargin=1.4em}
\newcommand{\btn}[1]{\textbf{#1}}
\newcommand{\tab}[1]{\textsf{\textbf{#1}}}
\newcommand{\tip}[1]{\par\smallskip\noindent\colorbox{tipc}{\parbox{\dimexpr\linewidth-2\fboxsep}{\textbf{Tip.} #1}}\par\smallskip}
\newcommand{\note}[1]{\par\smallskip\noindent\colorbox{soft}{\parbox{\dimexpr\linewidth-2\fboxsep}{\textbf{Note.} #1}}\par\smallskip}
\newcommand{\caution}[1]{\par\smallskip\noindent\colorbox{warn}{\parbox{\dimexpr\linewidth-2\fboxsep}{\textbf{Caution.} #1}}\par\smallskip}
\pagestyle{fancy}\fancyhf{}\rhead{\small DINO-4DSTEM user manual}\lhead{\small\leftmark}\rfoot{\small\thepage}
\renewcommand{\headrulewidth}{0.3pt}

\begin{document}
\begin{titlepage}
\vspace*{3cm}
\begin{center}
{\Huge\bfseries\color{deep} DINO-4DSTEM}\\[6pt]
{\LARGE A label-free tool for reading 4D-STEM diffraction maps}\\[10pt]
{\large User manual}\\[4pt]
{\normalsize For first-time and occasional users. No programming needed.}\\[40pt]
\colorbox{soft}{\parbox{0.85\linewidth}{\centering\vspace{4pt}
This program takes a 4D-STEM dataset and automatically sorts every scan position
into a small number of ``classes'' that share the same diffraction, giving you a
colour map of your sample, without labels, without choosing the number of classes,
and without setting peak thresholds. It can also run the classical comparison
methods (NMF, ACOM, SAM) and tell you what the classes mean, all in one window.
\vspace{4pt}}}
\vfill
{\small Screenshots will be added in a later revision; this edition is text-only.}
\end{center}
\end{titlepage}

\tableofcontents
\clearpage

% =====================================================================
\section{What this program does (in plain words)}
A 4D-STEM experiment scans a tiny electron beam across your sample and, at every
scan position, records a 2D \emph{diffraction pattern}. The result is a large
4-dimensional dataset: a 2D grid of positions, each holding a 2D diffraction
image. Looking through thousands of these patterns by hand is slow and subjective.

DINO-4DSTEM learns, on its own, to group scan positions whose diffraction looks
alike. The output is a \emph{class map}: an image of your scan where each colour is
one group (``class'') of similar diffraction. From there you can look at the
average diffraction of each class, ask the program what physically distinguishes
the classes, and compare the result against the traditional methods.

\textbf{Three things make it easy to use:}
\begin{itemize}[itemsep=2pt]
\item \textbf{No labelling.} You do not have to tell it what to look for.
\item \textbf{No ``number of classes'' to guess.} It finds how many are needed.
\item \textbf{No thresholds to tune.} It does not rely on peak-finding settings
that you would otherwise have to adjust by trial and error.
\end{itemize}

\note{You can use the program in two ways. (1) \emph{Train} a fresh model on your
own data (takes a while, uses the graphics card), or (2) \emph{load an existing
trained model} and just look at and analyse the results (fast). Both are covered
below.}

% =====================================================================
\section{Starting the program and the main window}
\subsection{Starting it}
Launch the program the way it was set up on your computer (a desktop shortcut, or
by running \texttt{gui\_dino4dstem.py} in the project's Python environment). A
single window titled \tab{DINO-4DSTEM} opens. A small black text window may also
appear behind it; that is the program's log. You can ignore it, but do not close
it, as closing it closes the program.

\tip{Maximise the window (the square button at the top-right) before you start, so
all the controls and images are visible.}

\subsection{The top bar (always visible)}
Across the very top of the window, from left to right:
\begin{itemize}[itemsep=3pt]
\item \btn{Inspect runs (HTML)} : opens a summary page of all the models you have
trained, in your web browser. Useful for finding an earlier result.
\item \btn{Assistant} : opens a plain-English chat helper that can carry out
common actions for you by typing requests in everyday language. (Optional; may not
be present in every version.)
\item \textbf{dataset badge} (``dataset: click to load a cube''): shows which data
file is currently loaded. Click it to load one.
\item \textbf{run badge} (``run: none''): shows which trained model is currently
in use. Click it to pick one.
\item \textbf{inference dot} : a small status light that tells you whether a class
map has been computed yet.
\item \textbf{real space} (\,nm/px\,): the distance between neighbouring scan
positions, i.e.\ your scan step size in nanometres. Used only for the scale bar and
axis labels on real-space images.
\item \textbf{reciprocal} (\,nm$^{-1}$/px\,): the calibration of the detector, in
nm$^{-1}$ per original detector pixel, as reported by your microscope for the
camera length you used. Used for the diffraction scale bars and the 1D plots.
\end{itemize}
\note{You only ever type these two calibration numbers once. Every plot in the
program automatically applies the right scaling, so you do not need to recompute
anything when the program resizes images internally.}

\subsection{How the tabs are organised}
The work is grouped into five big tabs along the top:
\begin{itemize}[itemsep=2pt]
\item \tab{Data} : load your 4D-STEM file and prepare it.
\item \tab{Model} : train a model, watch it learn, and apply it to other data.
\item \tab{Analysis} : look at the class map and understand what the classes mean.
\item \tab{Clustering} : the classical comparison methods (NMF, pretrained-DINO, SAM).
\item \tab{Diffraction} : per-pattern tools (peak/blob detection and crystal
orientation mapping, ACOM).
\end{itemize}
Each big tab contains smaller sub-tabs (for example \tab{Data} contains
\tab{Load + Pre-process} and \tab{Synthetic}). Click a big tab first, then the
sub-tab you want.

% =====================================================================
\section{The five-minute workflow}
If you just want a class map of your data, here is the whole path. Each step is
explained in detail later.
\begin{enumerate}[itemsep=3pt]
\item \tab{Data \textrightarrow{} Load + Pre-process}: click \btn{Browse}, choose
your 4D-STEM file, click \btn{Load}. Set \textbf{vmax} so the rings are visible
without the centre being a white blob. Leave the crop/mask at their defaults
unless you have a reason to change them.
\item Enter your \textbf{real space} and \textbf{reciprocal} calibration in the top
bar.
\item \tab{Model \textrightarrow{} Training}: pick your \textbf{Sample}, leave
\textbf{Use paper recipe (locked)} on, and click \btn{Train}. Wait for it to
finish. \emph{Or}, to skip training, click the \textbf{run badge} and load an
existing trained model.
\item \tab{Model \textrightarrow{} Eval} (or \tab{Analysis \textrightarrow{}
Post-hoc}): view the class map.
\item \tab{Analysis \textrightarrow{} Post-hoc \textrightarrow{} Interpretation}:
click \btn{Run interpretation} to get figures and a written report on what the
classes mean.
\end{enumerate}

% =====================================================================
\section{Data: loading and preparing your file}
\tab{Data \textrightarrow{} Pre-processing \textrightarrow{} Load + Pre-process}

This is where you open your data and get it ready. It is the first tab you use.

\subsection{Loading a file}
\begin{itemize}[itemsep=2pt]
\item \btn{Browse} then \btn{Load}: open a single 4D-STEM cube. Supported file
types include \texttt{.prz}, \texttt{.npy}, \texttt{.npz}, and \texttt{.h5}.
\item \btn{Browse multi\dots} : open several files together (for example four
quadrants of one acquisition).
\end{itemize}
After loading, the left side shows one diffraction pattern, and the bottom says how
the data is shaped. If nothing appears, check the log window for a message.

\subsection{Looking through the patterns}
\begin{itemize}[itemsep=2pt]
\item \textbf{Pattern index} slider: moves through the scan positions. Click the
slider or the image first, then use the left/right arrow keys for single steps, or
Page-Up/Page-Down for jumps of ten. \btn{Random} jumps to a random position.
\item \textbf{vmax}: the brightness ceiling. Lower it until you can see the faint
diffraction rings; raise it if the image is too dark. This only changes the
display and how the model sees intensities, not your raw file.
\item \textbf{cmap}: the colour scheme of the diffraction image.
\item \textbf{log-stretch}: compresses very bright and very faint together; helpful
when the central beam is far brighter than the rings.
\end{itemize}
\tip{Choosing a good \textbf{vmax} is the single most useful thing you do here. Aim
for a picture where both the bright centre and the outer rings are visible.}

\subsection{Crop and central-beam mask}
These define what part of each diffraction pattern the model looks at.
\begin{itemize}[itemsep=2pt]
\item \textbf{center crop (px, in 192 frame)}: how far out from the centre the
model looks. Keep your rings of interest comfortably inside this circle.
\item \textbf{polar mask cols (low-r mask in polar)}: hides the very centre (the
direct beam) so the intense central spot does not dominate. Too small lets the beam
bleed in; too large eats the first ring.
\end{itemize}
\note{The defaults are sensible for most molecular-film data. Only change them if
you know the ring you care about is being cut off (too-small crop) or the centre is
swamping everything (too-small mask).}

\subsection{Other preparation tools (optional)}
\begin{itemize}[itemsep=3pt]
\item \textbf{Binning (down-sample)}: \textbf{bin factor n} then
\btn{Apply binning} writes a new, smaller, faster file. Use only if loading is slow.
Choose \textbf{real-space} binning to merge scan positions or \textbf{q-space} to
merge detector pixels.
\item \textbf{Gaussian blur}: \textbf{$\sigma$ (px)} then \btn{Apply blur to
dataset} writes a gently smoothed copy. Reduces speckle noise on very low-dose data.
\item \textbf{Ellipticity correction}: fixes a slightly oval detector. Use
\btn{Auto-fit}, or type \textbf{a/b} and \textbf{$\theta$} and click \btn{Apply}.
\btn{Reset} returns to no correction.
\item \textbf{NBED centering}: \btn{COM-center the direct beam} /
\btn{Center entire dataset} re-centres each pattern on the beam if your data drifts.
\item \textbf{Pair labelling (pre-train)}: an advanced option to give the model a
few ``these two look the same / different'' hints. Most users never need this.
\end{itemize}
\caution{The buttons that say ``write new .cube.npy'' create a new file on disk;
they do not overwrite your original. Note the new file name so you load the right
one for training.}

\subsection{Synthetic (advanced)}
\tab{Data \textrightarrow{} Synthetic} simulates diffraction data from a known
crystal structure, to test the method. It is for advanced users and is not needed
for normal analysis.

% =====================================================================
\section{Model: training and applying a model}
\subsection{Training}
\tab{Model \textrightarrow{} Training \textrightarrow{} Training}

This teaches a model on your loaded data.
\begin{enumerate}[itemsep=2pt]
\item Choose the \textbf{Sample} (the data you prepared).
\item Optionally type a \textbf{Run name} so you can find the result later.
\item Leave \btn{Use paper recipe (locked)} switched on. This applies the
validated default settings, so you do not have to choose anything.
\item Click \btn{Train}. Progress shows in the status line; you can watch the map
form live on the \tab{Eval} sub-tab. Click \btn{Stop} to end early.
\end{enumerate}
\begin{itemize}[itemsep=2pt]
\item \btn{Refresh crop/mask/COM from PrePanel}: copies the crop and mask you set
on the Data tab into the training settings, so the model trains on exactly what you
previewed.
\item To change advanced settings, click the lock switch to unlock the recipe, then
\btn{Apply changes}. \btn{Apply auto-default ($\lambda$\_pair = 0.1)} turns on the
optional pair-hint term if you provided pair labels.
\end{itemize}
\caution{Training uses the graphics card and can take a while. Avoid running other
heavy GPU work at the same time. If your data loading feels stuck, that is usually
disk speed, not a crash; give it time.}

\subsection{Eval: viewing the class map}
\tab{Model \textrightarrow{} Training \textrightarrow{} Eval}
\begin{itemize}[itemsep=2pt]
\item During training this updates automatically (``awaiting first
checkpoint\dots'' then the live map).
\item To view a finished model: \btn{Load run dir\dots{} (model)}, pick the run
folder, and click \btn{Render}. \btn{Browse cube\dots} lets you apply the model to
a different data file.
\end{itemize}

\subsection{Transfer: applying one model to many files}
\tab{Model \textrightarrow{} Transfer}

Use a trained model on a whole folder of other datasets in one go.
\begin{enumerate}[itemsep=2pt]
\item Point \textbf{Cubes folder} at the folder, click \btn{Refresh}, and tick the
files you want (up to twenty).
\item Pick the starting model under \textbf{ckpt}.
\item Click \btn{Run sequential transfer}. \btn{Stop} ends it.
\end{enumerate}
\begin{itemize}[itemsep=2pt]
\item \textbf{freeze encoder (head-only adapt)}: a faster, lighter adaptation.
\item \textbf{Chain ckpts}: each file starts from the previous file's result;
useful for a time series.
\end{itemize}

% =====================================================================
\section{Analysis: understanding the class map}
\subsection{Post-hoc: explore the map}
\tab{Analysis \textrightarrow{} Post-hoc \textrightarrow{} Analysis}

This is the main place to look at results.
\begin{itemize}[itemsep=3pt]
\item \btn{Load run dir\dots} and \btn{Browse cube\dots}: load a trained model and
its data (or use the ones already in the session).
\item \textbf{Click the class map} to inspect a position: it shows that position's
raw pattern. Clicking a class shows the \textbf{class average} diffraction
(``class avg (top 200)''), which is much cleaner than a single noisy frame.
\item \textbf{Attribution (Avg + GradCAM + IG)}: \btn{Compute GradCAM + IG} makes a
picture showing \emph{which part of the diffraction the model used} to decide a
class. Choose the \textbf{CAM layer} if prompted; the defaults are fine.
\item \textbf{Merge classes}: tick two or more classes and click \btn{Apply merge}
to combine ones you consider the same. \btn{Reset} undoes it.
\item \textbf{Fine-tune with labels}: an advanced option to nudge the model using a
few of your own pair labels.
\item \btn{Send grain to ACOM tab}: hands the selected region to the orientation-
mapping tool (see Diffraction).
\item \textbf{Saving}: \btn{Save snapshot}, \btn{Save panel (PDF+PNG)}, and
\btn{Report All (PPTX + PNGs)} export your figures.
\end{itemize}
\tip{The class average is your friend. A single low-dose pattern looks like noise;
the class average reveals the real rings and spots.}

\subsection{Interpretation: what do the classes mean?}
\tab{Analysis \textrightarrow{} Post-hoc \textrightarrow{} Interpretation}

This runs a battery of checks and writes a short report.
\begin{enumerate}[itemsep=2pt]
\item Tick the \textbf{Analyses to run} you want (leave all on for a full report).
\item Click \btn{Run interpretation}. When done, use the \textbf{Figure} selector
to flip through the results, or \btn{Open output folder} to find the files.
\item \btn{Refresh from Post-hoc} re-reads the currently loaded run.
\end{enumerate}
The report answers, in figures and text: which physical quantities the classes
track (for example crystallinity or texture), whether the splits are real or
redundant, and how the model compares to the classical methods.

% =====================================================================
\section{Clustering: the classical comparison methods}
These tabs let you run the traditional approaches on the \emph{same} data, so you
can compare them with the model directly.

\subsection{NMF + cluster}
\tab{Clustering \textrightarrow{} NMF + cluster}

Non-negative matrix factorization followed by clustering (a standard baseline).
\begin{enumerate}[itemsep=2pt]
\item \btn{Load .prz / .npy} (or use the session sample).
\item Set \textbf{n\_comp} (how many components) and pick the \textbf{method}
(\textbf{K-means} is the default; \textbf{Agglomerative}, \textbf{FCM},
\textbf{HDBSCAN} are alternatives). \textbf{K} is the number of clusters, or choose
\textbf{auto}.
\item Click \btn{Run}. \btn{Show class averages} and \btn{Open interactive map}
display the result; \btn{Save PNG} / \btn{Report All} export it.
\end{enumerate}

\subsection{DINO + cluster}
\tab{Clustering \textrightarrow{} DINO + cluster}

Uses a general, pre-trained image model (downloaded automatically) instead of one
trained on your data, then clusters its features. A second baseline. Controls
mirror the NMF tab, plus \btn{Show UMAP} (a 2D scatter of the patterns) and
\btn{Class attentions (grid)}.

\subsection{SAM}
\tab{Clustering \textrightarrow{} SAM}

Uses the ``Segment Anything'' model to outline features in a class-average
diffraction pattern.
\begin{itemize}[itemsep=2pt]
\item Set \textbf{SAM ckpt} to the model file, pick a \textbf{class}, and click
\btn{Run on this average}.
\item \btn{Apply to whole class} / \btn{Apply to ALL classes} repeat it broadly.
\item \btn{Save tuning to per-class config} remembers your settings per class.
\end{itemize}

% =====================================================================
\section{Diffraction: per-pattern tools}
\subsection{Blob (peak detection)}
\tab{Diffraction \textrightarrow{} Blob}

Finds spots/peaks in a diffraction image.
\begin{itemize}[itemsep=2pt]
\item Pick a \textbf{class} and click \btn{Run on current image}, or set a
\textbf{frame idx} / \textbf{grain (y, x)} to target a specific pattern.
\item \textbf{Method} chooses the detector; \textbf{log1p stretch} helps faint
peaks show up. \btn{Apply to whole class} / \btn{Apply to ALL classes} run broadly.
\end{itemize}

\subsection{ACOM (crystal orientation mapping)}
\tab{Diffraction \textrightarrow{} ACOM}

Matches each pattern to a known crystal structure (from a CIF file) to find crystal
\emph{orientation}. It works in clear steps, left to right:
\begin{enumerate}[itemsep=3pt]
\item \btn{Load source}: pulls the current data and class map from the session (no
separate load button; it uses the top-bar badges).
\item \btn{+ CIF\dots} then \btn{Build all}: add your structure file(s) and build
the reference reflections.
\item \btn{Detect peaks}: finds peaks in the patterns. \textbf{log stretch before
detection}, \textbf{Top-N}, and \textbf{skip if < N peaks} control sensitivity.
\item \textbf{Calibration}: set \textbf{1/\AA{} per px} (or click \btn{Auto-fit})
and \textbf{k\_max}.
\item \btn{Match single}: solve one position (type a scan (y, x) or click the class
map). \btn{Single-phase full} / \btn{Multi-phase full} solve the whole scan
(\textbf{full-dataset stride} trades speed for completeness; \textbf{corr
threshold} and \textbf{winner margin} control acceptance).
\item \btn{Overlay} shows the orientation map; tick \textbf{also show in-plane
orientation} for the rotation angle.
\end{enumerate}
\note{ACOM needs a structure file (CIF) and a few detection settings, which is
exactly the kind of manual tuning the model avoids. Use it as a cross-check or when
you specifically want crystal orientation, which the model deliberately ignores.}

% =====================================================================
\section{Where your results are saved}
\begin{itemize}[itemsep=2pt]
\item Trained models live in a \texttt{runs} folder, one folder per model, each
holding the model file, a settings file, a training log, and an \texttt{eval}
folder with the class map.
\item Figures and reports you export (snapshots, panels, ``Report All'') are
written next to the run or to a folder the button names; use \btn{Open output
folder} on the Interpretation tab to find them.
\item New data files you create (binned, blurred, centred) are written as new
\texttt{.cube.npy} files in the same place as the original.
\end{itemize}

% =====================================================================
\section{Common tasks at a glance}
\begin{itemize}[itemsep=4pt]
\item \textbf{Just show me a class map of my data, using an existing model.}
Click the \textbf{run badge} (top bar) and load the model; click the
\textbf{dataset badge} and load your file; go to \tab{Model \textrightarrow{} Eval}
and click \btn{Render}.
\item \textbf{Train a new model on my data.}
\tab{Data} load + set vmax; \tab{Model \textrightarrow{} Training}, pick the sample,
keep the locked recipe, \btn{Train}.
\item \textbf{Understand what the classes are.}
\tab{Analysis \textrightarrow{} Post-hoc \textrightarrow{} Interpretation},
\btn{Run interpretation}.
\item \textbf{Compare to the classical method.}
\tab{Clustering \textrightarrow{} NMF + cluster}, \btn{Run}, then compare its map to
the model's.
\item \textbf{Get crystal orientations.}
\tab{Diffraction \textrightarrow{} ACOM}, add a CIF, \btn{Build all},
\btn{Detect peaks}, \btn{Single-phase full}, \btn{Overlay}.
\item \textbf{Apply one model to many files.}
\tab{Model \textrightarrow{} Transfer}, pick the folder, tick files,
\btn{Run sequential transfer}.
\end{itemize}

% =====================================================================
\section{Troubleshooting}
\begin{itemize}[itemsep=3pt]
\item \textbf{Nothing happens after I click Load.} Check the small black log
window for a message; the file type or scan shape may not have been recognised.
Try \btn{Browse} again and confirm the file.
\item \textbf{The diffraction image is all white in the middle, or all black.}
Adjust \textbf{vmax} on the Data tab and turn on \textbf{log-stretch}.
\item \textbf{The class map looks empty or a single colour.} Make sure a model is
loaded (run badge) and a class map has been computed (the inference light); on the
Eval tab click \btn{Render}.
\item \textbf{Training seems frozen.} Low-dose data loads slowly; wait a few
minutes. If the graphics card is busy with other work, training will be slow.
\item \textbf{A figure looks cramped or mislabelled.} Resize the window larger and
re-render; many panels redraw to fit.
\item \textbf{I changed crop/mask but training ignored it.} On the Training tab
click \btn{Refresh crop/mask/COM from PrePanel} before \btn{Train}.
\end{itemize}

% =====================================================================
\section{Glossary}
\begin{description}[style=nextline,leftmargin=1.4em,itemsep=3pt]
\item[4D-STEM cube] The whole dataset: a 2D grid of scan positions, each with a 2D
diffraction pattern.
\item[Diffraction pattern] The 2D image recorded at one scan position; rings and
spots in it encode the local structure.
\item[Class / class map] A group of scan positions with similar diffraction, and
the colour image showing where each group sits on the sample.
\item[Class average] The averaged diffraction of all positions in a class; cleaner
than any single pattern.
\item[Training] Letting the model learn groupings from your data. Done once per
dataset; afterwards the model can be reloaded instantly.
\item[Run] One trained model and its outputs, stored in its own folder.
\item[vmax] The brightness ceiling used when showing and reading diffraction
intensities.
\item[Crop / mask] How much of each pattern the model looks at, and how much of the
bright central beam is hidden.
\item[NMF, SAM, ACOM] Classical analysis methods provided for comparison:
decomposition-plus-clustering, feature segmentation, and crystal-orientation
mapping, respectively.
\item[Calibration (real / reciprocal)] The two numbers, scan step size (nm/px) and
detector calibration (nm$^{-1}$/px), that set the scale bars and axes.
\end{description}

\vfill
\noindent\colorbox{soft}{\parbox{\dimexpr\linewidth-2\fboxsep}{\small
\textbf{Getting help.} If you are unsure which setting to use for a new material,
the program's defaults (the locked ``paper recipe'') are a safe starting point. The
on-screen ``\,?\,'' help buttons next to many controls explain each one in place.}}
\end{document}
